% ===== PRÉAMBULE CANONIQUE — à coller VERBATIM en tête de CHAQUE .tex =====
\documentclass[11pt,a4paper]{article}
\usepackage[utf8]{inputenc}\usepackage[T1]{fontenc}\usepackage{lmodern}
\usepackage[french]{babel}
\usepackage{amsmath,amssymb,mathtools}\usepackage{icomma}
\usepackage[version=4]{mhchem}          % \ce{...} : équations & formules
\usepackage{siunitx}\sisetup{locale=FR} % \qty{}{}, \si{}, \ang{}
\usepackage{chemfig}                    % \chemfig{...} : structures développées
\usepackage{xcolor}\usepackage{colortbl}
\usepackage{tikz}\usepackage{pgfplots}\pgfplotsset{compat=1.18}
\usepackage[most]{tcolorbox}\usepackage{enumitem}\usepackage{array,booktabs}
\usepackage[a4paper,margin=2.2cm]{geometry}
\usetikzlibrary{arrows.meta,calc,angles,quotes,positioning,patterns,backgrounds,shapes.geometric}
\definecolor{primary}{RGB}{222,49,99}\definecolor{primarydark}{RGB}{150,22,55}
\definecolor{defcol}{RGB}{0,114,178}\definecolor{methcol}{RGB}{52,92,148}
\definecolor{excol}{RGB}{0,135,90}\definecolor{attncol}{RGB}{200,40,55}
\tcbset{boxbase/.style={enhanced,breakable,boxrule=0.9pt,arc=2.5pt,
  left=3mm,right=3mm,top=2.3mm,bottom=2.3mm,fonttitle=\bfseries,coltitle=white}}
% --- boîtes TUTORIEL (noms EXACTS, ne pas renommer) ---
\newtcolorbox{cle}[1][]{boxbase,colback=defcol!6,colframe=defcol,
  title={\textsc{À retenir}\ifx\relax#1\relax\relax\else\textmd{ : \itshape #1}\fi}}
\newtcolorbox{methode}[1][]{boxbase,colback=methcol!7,colframe=methcol,sharp corners,
  title={\textsc{Méthode}\ifx\relax#1\relax\relax\else\textmd{ --- \itshape #1}\fi}}
\newtcolorbox{exemplebox}[1][]{boxbase,colback=excol!5,colframe=excol,
  title={\textsc{Exemple}\ifx\relax#1\relax\relax\else\textmd{ : \itshape #1}\fi}}
\newtcolorbox{piege}[1][]{boxbase,colback=attncol!6,colframe=attncol,
  title={\textsc{Piège}\ifx\relax#1\relax\relax\else\textmd{ : \itshape #1}\fi}}
\newtcolorbox{remarque}[1][]{boxbase,colback=black!4,colframe=black!45,
  title={\textsc{Remarque}\ifx\relax#1\relax\relax\else\textmd{ : \itshape #1}\fi}}
% --- boîtes EXERCICE / CORRIGÉ (noms EXACTS, auto-comptées) ---
\newtcolorbox[auto counter]{exo}[1][]{boxbase,colback=primary!4,colframe=primary,
  title={\textsc{Exercice}~\thetcbcounter\ifx\relax#1\relax\relax\else\textmd{ --- \itshape #1}\fi}}
\newtcolorbox[auto counter]{corrige}[1][]{boxbase,colback=excol!4,colframe=excol,
  title={\textsc{Corrigé}~\thetcbcounter\ifx\relax#1\relax\relax\else\textmd{ --- \itshape #1}\fi}}
\newcommand{\R}{\mathbb{R}}\newcommand{\N}{\mathbb{N}}\newcommand{\Z}{\mathbb{Z}}\newcommand{\ds}{\displaystyle}
% (un \usepackage{fancyhdr}+en-tête est possible ; il est ignoré côté web)
% =========================================================================
\begin{document}

\begin{corrige}[n.o.\ dans un ion polyatomique]
On pose $x$ = n.o.\ de l'atome central, $\ce{O}=-\mathrm{II}$, et somme $=$ charge de l'ion.
\begin{enumerate}[label=\alph*)]
  \item \ce{SO4^2-} : $x+4(-2)=-2 \Rightarrow x=+\mathrm{VI}$.
  \item \ce{NO3-} : $x+3(-2)=-1 \Rightarrow x=+\mathrm{V}$.
  \item \ce{ClO4-} : $x+4(-2)=-1 \Rightarrow x=+\mathrm{VII}$.
  \item \ce{PO4^3-} : $x+4(-2)=-3 \Rightarrow x=+\mathrm{V}$.
\end{enumerate}
\end{corrige}

\begin{corrige}[quand l'élément central est présent deux fois]
On n'oublie pas le coefficient $2$ devant $x$.
\begin{enumerate}[label=\alph*)]
  \item \ce{Cr2O7^2-} : $2x+7(-2)=-2 \Rightarrow 2x=+12 \Rightarrow x=+\mathrm{VI}$.
  \item \ce{CrO4^2-} : $x+4(-2)=-2 \Rightarrow x=+\mathrm{VI}$ (même degré que le dichromate).
  \item \ce{S2O3^2-} : $2x+3(-2)=-2 \Rightarrow 2x=+4 \Rightarrow x=+\mathrm{II}$ (n.o.\ \emph{moyen} du soufre).
\end{enumerate}
\end{corrige}

\begin{corrige}[un même élément, plusieurs n.o.]
$\ce{O}=-\mathrm{II}$ à chaque fois ; on ajuste selon la charge globale.
\begin{enumerate}[label=\alph*)]
  \item \ce{MnO2} (neutre) : $x+2(-2)=0 \Rightarrow x=+\mathrm{IV}$.
  \item \ce{MnO4-} : $x+4(-2)=-1 \Rightarrow x=+\mathrm{VII}$.
  \item \ce{MnO4^2-} : $x+4(-2)=-2 \Rightarrow x=+\mathrm{VI}$.
\end{enumerate}
\end{corrige}

\begin{corrige}[n.o.\ du métal à valence variable dans un oxyde]
Molécules neutres, $\ce{O}=-\mathrm{II}$.
\begin{enumerate}[label=\alph*)]
  \item \ce{Cu2O} : $2x+(-2)=0 \Rightarrow 2x=+2 \Rightarrow x=+\mathrm{I}$ (cuivre(I)).
  \item \ce{Fe2O3} : $2x+3(-2)=0 \Rightarrow 2x=+6 \Rightarrow x=+\mathrm{III}$ (fer(III)).
  \item \ce{SnO2} : $x+2(-2)=0 \Rightarrow x=+\mathrm{IV}$ (étain(IV)).
\end{enumerate}
\end{corrige}

\begin{corrige}[n.o.\ dans un oxacide neutre]
On combine $\ce{H}=+\mathrm{I}$, $\ce{O}=-\mathrm{II}$, molécule neutre. On note $x$ le n.o.\ central.
\begin{enumerate}[label=\alph*)]
  \item \ce{H2SO4} : $2(+1)+x+4(-2)=0 \Rightarrow x-6=0 \Rightarrow x=+\mathrm{VI}$.
  \item \ce{HNO2} : $(+1)+x+2(-2)=0 \Rightarrow x-3=0 \Rightarrow x=+\mathrm{III}$.
  \item \ce{H3PO4} : $3(+1)+x+4(-2)=0 \Rightarrow x-5=0 \Rightarrow x=+\mathrm{V}$.
\end{enumerate}
\end{corrige}

\end{document}
